\documentclass[a4paper,11pt]{ctexart}
\usepackage{amsmath, amssymb, amsfonts}
\usepackage{geometry}
\geometry{left=1.5cm, right=1.5cm, top=2.0cm, bottom=2.0cm, headsep=0.3cm, footskip=1cm}
\usepackage{listings}
\usepackage{xcolor}
\usepackage{tikz}
\usepackage{pgfplots}
\usepackage{hyperref}
\usepackage{fancyhdr}
\usepackage{tcolorbox}
\tcbuselibrary{breakable, skins}
\usepackage{array}
\usepackage{booktabs}
\usepackage[shortlabels]{enumitem}
\usepackage{needspace}
\usepackage{multirow}

\AtBeginDocument{
  \hypersetup{colorlinks=true, linkcolor=blue, filecolor=magenta, urlcolor=cyan}
}

\setlist{nosep, itemsep=0.8ex, parsep=0.1em}
\setlist[itemize]{leftmargin=1.8em}
\setlist[enumerate]{leftmargin=1.8em}

% ===== 颜色定义 =====
\definecolor{codegreen}{rgb}{0,0.6,0}
\definecolor{codegray}{rgb}{0.5,0.5,0.5}
\definecolor{codepurple}{rgb}{0.58,0,0.82}
\definecolor{codebg}{rgb}{0.98,0.98,0.98}
\definecolor{tipsblue}{rgb}{0.85,0.93,1.0}
\definecolor{highlightyellow}{rgb}{1.0,0.97,0.8}

% ===== 代码块样式 =====
\lstdefinestyle{mystyle}{
    backgroundcolor=\color{codebg},
    commentstyle=\color{codegreen},
    keywordstyle=\color{magenta}\bfseries,
    numberstyle=\tiny\color{codegray},
    stringstyle=\color{codepurple},
    basicstyle=\ttfamily\small,
    breakatwhitespace=false,
    breaklines=true,
    captionpos=b,
    keepspaces=true,
    numbers=left,
    numbersep=6pt,
    showspaces=false,
    showstringspaces=false,
    showtabs=false,
    tabsize=2,
    language=Python,
    frame=single,
    framerule=0.5pt,
    rulecolor=\color{gray!40}
}
\lstset{style=mystyle}

% ===== 彩色框 =====
\newtcolorbox{problembox}[1][]{
    colback=white,
    colframe=blue!60!black,
    title=\textbf{#1},
    fonttitle=\bfseries\small,
    boxrule=1pt,
    arc=3pt, outer arc=3pt,
    coltitle=black,
    before skip=10pt, after skip=8pt
}

\newtcolorbox{solutionbox}[1][]{
    enhanced, breakable,
    colback=green!3!white,
    colframe=green!50!black,
    title=\textbf{#1},
    fonttitle=\bfseries\small,
    boxrule=1pt,
    arc=3pt, outer arc=3pt,
    coltitle=black,
    before skip=8pt, after skip=10pt
}

\newtcolorbox{metaphorbox}[1][]{
    colback=orange!5!white,
    colframe=orange!70!black,
    title=\textbf{#1},
    fonttitle=\bfseries\small,
    boxrule=1pt,
    arc=3pt, outer arc=3pt,
    coltitle=black,
    before skip=8pt, after skip=8pt
}

\newtcolorbox{beginnerbox}[1][]{
    colback=tipsblue,
    colframe=blue!50!black,
    title=\textbf{#1},
    fonttitle=\bfseries\small,
    boxrule=1pt,
    arc=3pt, outer arc=3pt,
    coltitle=black,
    before skip=8pt, after skip=8pt
}

\newtcolorbox{appbox}[1][]{
    colback=green!3!white,
    colframe=green!60!black,
    title=\textbf{#1},
    fonttitle=\bfseries\small,
    boxrule=1pt,
    arc=3pt, outer arc=3pt,
    coltitle=black,
    before skip=8pt, after skip=10pt
}

\newtcolorbox{keybox}[1][]{
    colback=highlightyellow,
    colframe=orange!80!black,
    title=\textbf{#1},
    fonttitle=\bfseries\small,
    boxrule=1pt,
    arc=3pt, outer arc=3pt,
    coltitle=black,
    before skip=8pt, after skip=8pt
}

% ===== 页眉页脚 =====
\pagestyle{fancy}
\fancyhf{}
\fancyhead[R]{机器学习-第10章}
\fancyhead[L]{无监督学习 · 练习题}
\fancyfoot[C]{\thepage}
\addtolength{\headheight}{2pt}

\parskip=2pt plus 1pt minus 0.5pt
\linespread{1.35}
\raggedbottom
\widowpenalty=10000
\clubpenalty=10000

\title{\textbf{机器学习\\第10章\quad 无监督学习}\\[0.5em]
{\large\color{gray}——当数据没有标签，机器该怎样自己发现规律？}}
\date{\today}

\begin{document}

\maketitle

% ===== 章节导读 =====
\begin{beginnerbox}[本章题目导览：你将挑战哪些核心问题？]
    在监督学习中，每条数据都有一个正确答案（标签）告诉模型"这是猫""这是狗"。但现实中，大量数据根本没有标签——海量的用户行为、基因序列、图片……人工标注代价极高。\textbf{无监督学习}的任务就是：在没有标签的情况下，从数据中自动发现结构和规律。本章 7 道题覆盖聚类与降维两大核心方向：

    \begin{itemize}
        \item \textbf{Q10-01（概念辨析）}：有监督 vs 无监督——两种学习范式的根本区别，聚类到底在做什么？
        \item \textbf{Q10-02（算法推导）}：K-means 工作流程手动推导——给定数据点，一步步演算质心更新过程，理解收敛本质。
        \item \textbf{Q10-03（对比分析）}：三种聚类算法横向对比——K-means、层次聚类、密度聚类各有什么适用场景和缺陷？
        \item \textbf{Q10-04（代码实践）}：K-means API 使用——\texttt{sklearn} 的 \texttt{KMeans} 参数解析，如何用代码完成一次完整的聚类分析？
        \item \textbf{Q10-05（评估指标）}：聚类模型评估——轮廓系数、WCSS、CH 指数分别衡量什么？肘部法如何选 $K$？
        \item \textbf{Q10-06（数学理解）}：奇异值分解（SVD）——矩阵分解的核心原理，为什么 SVD 能实现降维？
        \item \textbf{Q10-07（综合应用）}：主成分分析（PCA）——PCA 和 SVD 的关系，用 SVD 实现 PCA 降维的完整流程是什么？
    \end{itemize}
\end{beginnerbox}

\section{第 10 章\quad 无监督学习：聚类与降维}

前面所有章节学的都是\textbf{监督学习}：每条训练数据都有一个"正确答案"作为老师，模型在老师的监督下学习如何从输入预测输出。

但现在换一种场景：

\begin{itemize}
    \item 你有 100 万条用户的购物记录，没有人告诉你用户属于哪类群体；
    \item 你有一张 $1024 \times 1024$ 的高清图片，想压缩成更小的尺寸但不失真太多；
    \item 你有 500 维的基因表达数据，想在二维平面上看清楚哪些样本相似。
\end{itemize}

这类问题没有标签，无法用交叉熵损失来训练。我们需要\textbf{无监督学习}：让模型自己在数据中发现结构。

\begin{metaphorbox}[先建立一个总感觉：无监督学习在做什么]
    \textbf{想象你是一个图书管理员，面前堆着一万本没有任何分类标签的书。}

    \begin{itemize}
        \item \textbf{聚类}：把书摆来摆去，按内容相似度自动归堆——武侠小说堆一起，技术书堆一起，虽然没人告诉你怎么分，但你凭书的"相似感"把它们分开了。
        \item \textbf{降维}：书太多放不下，你决定用一句话概括每本书的精华——从厚厚的书变成一句核心摘要，用更少的信息保留最重要的内容。
    \end{itemize}

    \textbf{无监督学习的核心}：没有老师打分，模型自己发现数据内部的\textbf{结构}和\textbf{规律}。
\end{metaphorbox}


%% ============================================================
\Needspace{5\baselineskip}
\begin{problembox}[Q10-01\quad 有监督 vs 无监督——范式对比与聚类本质]
    \textbf{题目描述：}\\
    理解无监督学习最重要的一步，是明白它和监督学习的根本区别，以及聚类这件事的本质目标。

    \vspace{0.3cm}
    \textbf{问题：}
    \begin{enumerate}[(1)]
        \item 填写下表，对比有监督学习与无监督学习的关键区别：

        \begin{center}
        \begin{tabular}{p{3.5cm} p{4.5cm} p{4.5cm}}
            \toprule
            对比维度 & 有监督学习 & 无监督学习 \\
            \midrule
            训练数据 & 带标签的样本 $\{(x_i, y_i)\}$ & \underline{\hspace{3.5cm}} \\[1.2em]
            学习目标 & 学习从输入到标签的映射 $f: x \to y$ & \underline{\hspace{3.5cm}} \\[1.2em]
            典型任务 & 分类、回归 & \underline{\hspace{3.5cm}} \\[1.2em]
            评估方式 & 对比预测值与真实标签（如准确率） & \underline{\hspace{3.5cm}} \\
            \bottomrule
        \end{tabular}
        \end{center}

        \item 聚类（Clustering）的核心目标用一句话描述是什么？它应满足哪两个基本原则（关于簇内和簇间的相似度）？

        \item 以下哪些任务属于无监督学习？（多选）
        \begin{enumerate}[A.]
            \item 给邮件自动打"垃圾邮件/正常邮件"标签（已有历史标注数据）
            \item 对电商平台的用户按行为特征自动分群，用于精准营销
            \item 将高维医疗影像数据压缩到二维平面以便可视化
            \item 根据历史房价数据预测未来房价
            \item 在没有类别标注的新闻语料库中，自动发现潜在的话题分布
        \end{enumerate}
    \end{enumerate}

    \vspace{0.2cm}
    {\color{gray}\small\textit{来源溯源：[文档第 10 章 10.1.1 聚类简介]}}
\end{problembox}

\begin{beginnerbox}[小白必读：为什么"发现结构"这么重要？]
    \begin{itemize}
        \item \textbf{标签太贵了}：ImageNet 有 120 万张图片，每张都需要人工标注，代价极高。但互联网上有数十亿张无标签图片，无监督学习就是利用这些"免费数据"的工具。
        \item \textbf{很多时候不知道该分几类}：监督分类时，类别是预先定义好的。但聚类不需要预设类别，它可以帮你\textbf{发现数据中自然存在的分组}，这在探索性分析中非常有价值。
        \item \textbf{降维是数据分析的基础工具}：高维数据（几百维）直接分析很困难，降到 2-3 维后，人眼就能直接看出数据的结构——聚不聚集、有没有异常点。
    \end{itemize}
\end{beginnerbox}

\begin{solutionbox}[参考答案与详细解析]
    \textbf{(1) 对比表格填写：}

    \begin{center}
    \begin{tabular}{p{3.5cm} p{4.5cm} p{4.5cm}}
        \toprule
        对比维度 & 有监督学习 & 无监督学习 \\
        \midrule
        训练数据 & 带标签的样本 $\{(x_i, y_i)\}$ & 仅有特征，无标签 $\{x_i\}$ \\[1.2em]
        学习目标 & 学习从输入到标签的映射 & 发现数据内在结构（分布、分组、主要成分） \\[1.2em]
        典型任务 & 分类、回归 & 聚类、降维、异常检测、密度估计 \\[1.2em]
        评估方式 & 对比预测值与真实标签 & 内部指标（轮廓系数等）或下游任务效果 \\
        \bottomrule
    \end{tabular}
    \end{center}

    \textbf{(2) 聚类的核心目标：}

    将数据集划分为若干个簇（Cluster），使得\textbf{同一簇内的样本尽可能相似（簇内距离小）}，\textbf{不同簇之间的样本尽可能不同（簇间距离大）}。

    \textbf{(3) 属于无监督学习的任务：B、C、E。}

    \begin{itemize}
        \item A 是有监督分类（有历史标注）；
        \item D 是有监督回归（有历史房价作为标签）；
        \item B 是聚类（无预定义分群标签）；
        \item C 是降维可视化（无标签，发现数据结构）；
        \item E 是主题模型（无标注，发现话题分布）。
    \end{itemize}
\end{solutionbox}


%% ============================================================
\Needspace{5\baselineskip}
\begin{problembox}[Q10-02\quad K-means 工作流程手动推导]
    \textbf{题目描述：}\\
    K-means 是最经典的聚类算法之一。理解它的关键，是能够手动推演一次完整的迭代过程——从初始化质心，到分配簇，到更新质心，再到收敛判断。

    \vspace{0.3cm}
    \textbf{题目数据：}\\
    平面上有 6 个数据点：$A(1,1),\ B(1,3),\ C(2,2),\ D(8,1),\ E(8,3),\ F(9,2)$。
    设 $K=2$，初始质心为 $\mu_1 = (1,1)$（即点 $A$），$\mu_2 = (8,1)$（即点 $D$）。
    距离度量使用欧氏距离：$d(x, y) = \sqrt{(x_1-y_1)^2 + (x_2-y_2)^2}$。

    \vspace{0.2cm}
    \textbf{问题：}
    \begin{enumerate}[(1)]
        \item \textbf{第一次迭代——分配步骤}：计算每个点到 $\mu_1, \mu_2$ 的欧氏距离，将每个点分配到距离最近的质心所在的簇。填写下表：

        \begin{center}
        \begin{tabular}{ccccl}
            \toprule
            数据点 & 到 $\mu_1=(1,1)$ 的距离 & 到 $\mu_2=(8,1)$ 的距离 & 分配到簇 \\
            \midrule
            $A(1,1)$ & \underline{\quad} & \underline{\quad} & \underline{\quad} \\
            $B(1,3)$ & \underline{\quad} & \underline{\quad} & \underline{\quad} \\
            $C(2,2)$ & \underline{\quad} & \underline{\quad} & \underline{\quad} \\
            $D(8,1)$ & \underline{\quad} & \underline{\quad} & \underline{\quad} \\
            $E(8,3)$ & \underline{\quad} & \underline{\quad} & \underline{\quad} \\
            $F(9,2)$ & \underline{\quad} & \underline{\quad} & \underline{\quad} \\
            \bottomrule
        \end{tabular}
        \end{center}

        \item \textbf{第一次迭代——更新步骤}：根据第(1)步的分配结果，重新计算每个簇的质心（取簇内所有点坐标的均值）。

        \item \textbf{第二次迭代}：用更新后的质心重新分配各点，再次更新质心，判断算法是否已经收敛（质心不再变化）。

        \item K-means 的收敛条件是什么？它一定能找到全局最优解吗？为什么？
    \end{enumerate}

    \vspace{0.2cm}
    {\color{gray}\small\textit{来源溯源：[文档第 10 章 10.1.2 K均值聚类工作流程]}}
\end{problembox}

\begin{beginnerbox}[小白必读：K-means 的三个核心步骤]
    \begin{itemize}
        \item \textbf{Step 1——初始化}：随机选 $K$ 个点作为初始质心（或用 K-means++ 策略）。
        \item \textbf{Step 2——分配（E步）}：把每个数据点分配给距它最近的质心，形成 $K$ 个簇。
        \item \textbf{Step 3——更新（M步）}：对每个簇，把簇内所有点的坐标取均值，得到新质心。
        \item \textbf{重复 Step 2 和 Step 3}，直到质心不再移动（或移动距离小于阈值）为止。
    \end{itemize}
    "K-means"名字的含义：$K$ 是簇的数量，"means"就是"均值"——新质心是簇内数据点的\textbf{均值}。
\end{beginnerbox}

\begin{metaphorbox}[生活比喻：K-means 就像"自动选座位"]
    假设班里来了一批新同学，老师想把大家分成 $K$ 组，但不知道怎么分：
    \begin{itemize}
        \item \textbf{初始化}：随机叫 $K$ 个同学站到教室各个角落，作为"组长"。
        \item \textbf{分配步骤}：每个剩余同学走向离自己最近的组长，加入他的小组。
        \item \textbf{更新步骤}：每个小组重新找一个"重心位置"，让站在那里的人到组内所有成员距离之和最小——这个重心就是新的质心（往往不是原来的组长了）。
        \item \textbf{收敛}：等到没有人需要换组，也没有质心需要移动，分组就稳定下来了。
    \end{itemize}
\end{metaphorbox}

\begin{solutionbox}[参考答案与详细解析]
    \textbf{(1) 第一次迭代——分配步骤：}

    计算各点到两个质心的距离：

    \begin{center}
    \begin{tabular}{cccl}
        \toprule
        数据点 & 到 $\mu_1=(1,1)$ & 到 $\mu_2=(8,1)$ & 分配到簇 \\
        \midrule
        $A(1,1)$ & $0$ & $7$ & 簇1 \\
        $B(1,3)$ & $2$ & $\sqrt{53}\approx7.28$ & 簇1 \\
        $C(2,2)$ & $\sqrt{2}\approx1.41$ & $\sqrt{37}\approx6.08$ & 簇1 \\
        $D(8,1)$ & $7$ & $0$ & 簇2 \\
        $E(8,3)$ & $\sqrt{53}\approx7.28$ & $2$ & 簇2 \\
        $F(9,2)$ & $\sqrt{65}\approx8.06$ & $\sqrt{2}\approx1.41$ & 簇2 \\
        \bottomrule
    \end{tabular}
    \end{center}

    \textbf{(2) 第一次迭代——更新步骤：}

    \begin{itemize}
        \item 簇1 包含 $A(1,1), B(1,3), C(2,2)$：
        \[
        \mu_1^{\text{new}} = \left(\frac{1+1+2}{3},\ \frac{1+3+2}{3}\right) = \left(\frac{4}{3},\ 2\right) \approx (1.33,\ 2)
        \]
        \item 簇2 包含 $D(8,1), E(8,3), F(9,2)$：
        \[
        \mu_2^{\text{new}} = \left(\frac{8+8+9}{3},\ \frac{1+3+2}{3}\right) = \left(\frac{25}{3},\ 2\right) \approx (8.33,\ 2)
        \]
    \end{itemize}

    \textbf{(3) 第二次迭代：}

    用新质心 $\mu_1=(1.33, 2)$，$\mu_2=(8.33, 2)$ 重新分配各点：

    由于六个点的位置分布非常明显（左侧三点 vs 右侧三点），重新计算距离后，分配结果与第一次完全一致——簇1 仍为 $\{A,B,C\}$，簇2 仍为 $\{D,E,F\}$。质心不再移动，\textbf{算法在第二次迭代后收敛}。

    \textbf{(4) 收敛条件与全局最优：}

    \begin{keybox}[K-means 的收敛与局限性]
        \begin{itemize}
            \item \textbf{收敛条件}：质心在前后两次迭代中不再发生变化（或变化量小于设定阈值 \texttt{tol}）。理论上 K-means 一定会在有限步内收敛，因为每次迭代都不会增加目标函数（WCSS）的值。
            \item \textbf{不保证全局最优}：K-means 是贪心算法，只保证收敛到\textbf{局部最优解}。结果对初始质心的选取非常敏感——不同的初始质心可能导致完全不同的聚类结果。
            \item \textbf{解决方案}：使用 \textbf{K-means++} 初始化策略（倾向于选择彼此较远的初始质心），或多次随机初始化后取最优结果（\texttt{n\_init} 参数）。
        \end{itemize}
    \end{keybox}
\end{solutionbox}


%% ============================================================
\Needspace{5\baselineskip}
\begin{problembox}[Q10-03\quad 三种聚类算法横向对比]
    \textbf{题目描述：}\\
    除 K-means 外，层次聚类和密度聚类（DBSCAN）也是常用的聚类方法。三种算法各有所长，理解它们的差异有助于在实际问题中选择合适的工具。

    \vspace{0.3cm}
    \textbf{问题：}
    \begin{enumerate}[(1)]
        \item 填写下表，完成三种聚类算法的对比：

        \begin{center}
        \begin{tabular}{p{2.2cm} p{3.8cm} p{3.8cm} p{2.6cm}}
            \toprule
            算法 & 核心思想 & 主要局限性 & 需要预设 $K$？ \\
            \midrule
            K-means & 迭代更新质心，最小化\underline{\hspace{1.5cm}} & 对\underline{\hspace{1.5cm}}敏感，只能发现\underline{\hspace{1.5cm}}形状的簇 & \underline{\quad} \\[2em]
            层次聚类 & 自底向上合并（凝聚型）或自顶向下分裂，形成\underline{\hspace{1.5cm}} & 计算复杂度高：$O(\underline{\quad})$，无法处理超大规模数据 & \underline{\quad} \\[2em]
            密度聚类 (DBSCAN) & 寻找被低密度区域分隔的高密度区域，将密度相连的点归为一簇 & 对\underline{\hspace{1.5cm}}参数敏感，高维数据效果差 & \underline{\quad} \\
            \bottomrule
        \end{tabular}
        \end{center}

        \item DBSCAN 中有两个核心参数：$\varepsilon$（邻域半径）和 \texttt{MinPts}（核心点的最小邻居数）。请解释：
        \begin{enumerate}[a.]
            \item 什么是"核心点"、"边界点"和"噪声点"？
            \item 如果把 $\varepsilon$ 调得很大，会发生什么？如果 $\varepsilon$ 很小呢？
        \end{enumerate}

        \item 下面四种数据分布，分别最适合用哪种聚类算法？请匹配并说明理由。
        \begin{enumerate}[A.]
            \item 三个球形簇，各簇大小相近，维度较低（如 2D）
            \item 两个月牙形（非凸）簇，中间无明显间隔
            \item 希望得到一棵完整的"树形层次结构"，可视化数据的多层级关系
            \item 数据中含有大量异常点（噪声），簇的形状不规则
        \end{enumerate}
    \end{enumerate}

    \vspace{0.2cm}
    {\color{gray}\small\textit{来源溯源：[文档第 10 章 10.1.2 层次聚类 / 密度聚类]}}
\end{problembox}

\begin{beginnerbox}[小白必读：为什么要有这么多种聚类算法？]
    \begin{itemize}
        \item \textbf{K-means 的盲区}：K-means 假设每个簇是"球形的"，且大小相近。面对月牙形、环形等不规则簇，它会彻底失效——因为质心（均值点）根本无法代表非凸形状的中心。
        \item \textbf{层次聚类的优势}：不需要预设 $K$，生成的树状图（Dendrogram）保留了数据的完整层次关系，可以在任意层"切割"得到不同粒度的聚类结果。
        \item \textbf{DBSCAN 的独特能力}：天然支持发现任意形状的簇，并能自动识别噪声点（标记为 $-1$）。这对有噪声的真实数据非常宝贵——K-means 会把所有噪声点都强行分配到某个簇中。
    \end{itemize}
\end{beginnerbox}

\begin{solutionbox}[参考答案与详细解析]
    \textbf{(1) 对比表格填写：}

    \begin{center}
    \begin{tabular}{p{2.2cm} p{3.8cm} p{3.8cm} p{2.6cm}}
        \toprule
        算法 & 核心思想 & 主要局限性 & 需要预设 $K$？ \\
        \midrule
        K-means & 迭代更新质心，最小化簇内平方和（WCSS） & 对初始质心敏感，只能发现球形/凸形状的簇 & 是 \\[1em]
        层次聚类 & 自底向上合并，形成树状图（Dendrogram） & 复杂度 $O(n^2)$ 到 $O(n^3)$，无法处理超大规模数据 & 否 \\[1em]
        DBSCAN & 寻找高密度区域，密度相连点归一簇 & 对 $\varepsilon$、\texttt{MinPts} 参数敏感，高维数据效果差 & 否 \\
        \bottomrule
    \end{tabular}
    \end{center}

    \textbf{(2) DBSCAN 核心概念：}

    \begin{itemize}
        \item \textbf{核心点（Core Point）}：在半径 $\varepsilon$ 内邻居数 $\geq \texttt{MinPts}$ 的点，是高密度区域的"骨干"。
        \item \textbf{边界点（Border Point）}：邻居数 $< \texttt{MinPts}$，但在某个核心点的 $\varepsilon$ 邻域内，是簇的"边缘成员"。
        \item \textbf{噪声点（Noise Point）}：既不是核心点，也不在任何核心点的邻域内，标记为异常值（标签 $-1$）。
    \end{itemize}

    \begin{itemize}
        \item \textbf{$\varepsilon$ 很大}：越来越多的点相互"密度可达"，不同的簇可能被合并成一个大簇，极端情况下所有点都在同一个簇里。
        \item \textbf{$\varepsilon$ 很小}：大多数点找不到足够的邻居成为核心点，大量点被标记为噪声，簇数量极多且每个簇只含极少数点。
    \end{itemize}

    \textbf{(3) 场景匹配：}
    \begin{itemize}
        \item \textbf{A → K-means}：球形等大簇是 K-means 最适合的场景，计算效率高。
        \item \textbf{B → DBSCAN}：月牙形等非凸形状，K-means 无能为力，DBSCAN 能沿密度边界准确划分。
        \item \textbf{C → 层次聚类}：层次聚类生成完整的 Dendrogram，是可视化层次关系的首选。
        \item \textbf{D → DBSCAN}：DBSCAN 天然具备噪声识别能力，不规则形状也能处理。
    \end{itemize}
\end{solutionbox}


%% ============================================================
\Needspace{5\baselineskip}
\begin{problembox}[Q10-04\quad K-means API 使用：\texttt{sklearn.KMeans} 参数解析]
    \textbf{题目描述：}\\
    下面是一段使用 scikit-learn 的 \texttt{KMeans} 对鸢尾花数据集（只取前两个特征）进行聚类的完整代码，请结合代码回答问题。

    \begin{lstlisting}[language=Python]
from sklearn.cluster import KMeans
from sklearn.datasets import load_iris
from sklearn.preprocessing import StandardScaler
import matplotlib.pyplot as plt
import numpy as np

# 加载数据，只取前两个特征
iris = load_iris()
X = iris.data[:, :2]   # shape: (150, 2)

# 标准化
scaler = StandardScaler()
X_scaled = scaler.fit_transform(X)

# 训练 KMeans 模型
kmeans = KMeans(
    n_clusters=3,
    init='k-means++',
    n_init=10,
    max_iter=300,
    random_state=42
)
kmeans.fit(X_scaled)

# 获取结果
labels = kmeans.labels_           # 每个样本的簇标签
centers = kmeans.cluster_centers_ # 簇质心坐标
inertia = kmeans.inertia_         # WCSS 值
    \end{lstlisting}

    \vspace{0.2cm}
    \textbf{问题：}
    \begin{enumerate}[(1)]
        \item 解释 \texttt{KMeans} 构造函数中每个参数的含义：\texttt{n\_clusters}、\texttt{init}、\texttt{n\_init}、\texttt{max\_iter}、\texttt{random\_state}。

        \item 为什么在聚类之前要做\textbf{标准化}（\texttt{StandardScaler}）？如果不标准化，会发生什么问题？

        \item \texttt{kmeans.labels\_} 的形状和取值范围是什么？\texttt{kmeans.cluster\_centers\_} 的形状是什么？

        \item 训练完毕后，如果有一个新数据点 \texttt{x\_new = [[5.1, 3.5]]}，如何用训练好的模型预测它属于哪个簇？写出对应的代码。
    \end{enumerate}

    \vspace{0.2cm}
    {\color{gray}\small\textit{来源溯源：[文档第 10 章 10.1.3 K-means API 使用]}}
\end{problembox}

\begin{beginnerbox}[小白必读：为什么聚类一定要标准化？]
    \begin{itemize}
        \item \textbf{距离受量纲影响}：K-means 用欧氏距离衡量相似度。假设特征1的取值范围是 $[0, 1]$，特征2的取值范围是 $[0, 10000]$（如房价）。在计算距离时，特征2的差异会完全"碾压"特征1——相当于特征1根本没有参与聚类。
        \item \textbf{标准化消除量纲}：\texttt{StandardScaler} 将每个特征变换为均值为 0、标准差为 1 的分布，使不同量纲的特征在距离计算中"站在平等地位"上。
        \item \textbf{经验法则}：只要用到"距离"的算法（K-means、KNN、SVM、PCA……），都应该优先考虑是否需要标准化。
    \end{itemize}
\end{beginnerbox}

\begin{solutionbox}[参考答案与详细解析]
    \textbf{(1) 参数含义：}

    \begin{itemize}
        \item \textbf{\texttt{n\_clusters=3}}：聚成 $K=3$ 个簇，即最终输出 3 个类别。
        \item \textbf{\texttt{init='k-means++'}}：使用 K-means++ 策略初始化质心，倾向于将初始质心选得彼此尽量分散，减少陷入差结果的概率（相比纯随机初始化更稳定）。
        \item \textbf{\texttt{n\_init=10}}：用 10 种不同的随机初始质心各跑一次，取 WCSS 最小的结果作为最终输出，进一步降低局部最优风险。
        \item \textbf{\texttt{max\_iter=300}}：单次运行的最大迭代次数，防止在边缘情况下无限循环。
        \item \textbf{\texttt{random\_state=42}}：固定随机种子，保证结果可复现。
    \end{itemize}

    \textbf{(2) 标准化的必要性：}

    欧氏距离对各特征的量纲极其敏感。鸢尾花数据中，花萼长度（约 4--8 cm）和花萼宽度（约 2--5 cm）量纲相似，但若混入花瓣长度（0.1--7 cm）且未标准化，数值范围差异会导致距离主要由范围大的特征决定，范围小的特征几乎失效。标准化后，所有特征均等地参与距离计算。

    \textbf{(3) 结果形状：}

    \begin{itemize}
        \item \texttt{labels\_}：形状 \texttt{(150,)}，取值范围 $\{0, 1, 2\}$（即 $0$ 到 $K-1$），每个元素是对应样本被分配的簇编号。
        \item \texttt{cluster\_centers\_}：形状 \texttt{(3, 2)}，即 $K$ 个簇，每个簇的质心是一个 2 维向量（与输入特征维度一致）。
    \end{itemize}

    \textbf{(4) 预测新数据点的代码：}

    \begin{lstlisting}[language=Python]
x_new = [[5.1, 3.5]]
# 必须先用同一个 scaler 做标准化！
x_new_scaled = scaler.transform(x_new)
# 预测所属簇
pred_label = kmeans.predict(x_new_scaled)
print(f"新数据点被分配到簇: {pred_label[0]}")
    \end{lstlisting}

    \begin{keybox}[重要提醒：训练与预测时的标准化必须一致]
        训练时用 \texttt{fit\_transform(X)}，预测新数据时只能用 \texttt{transform(x\_new)}（不能再 \texttt{fit}）。原因：新数据必须用\textbf{训练集的均值和标准差}来缩放，才能和训练好的质心坐标处于同一空间中进行距离比较。
    \end{keybox}
\end{solutionbox}


%% ============================================================
\Needspace{5\baselineskip}
\begin{problembox}[Q10-05\quad 聚类模型评估：如何判断聚类效果好不好？]
    \textbf{题目描述：}\\
    聚类没有真实标签，无法用准确率评估。但有一系列\textbf{内部评估指标}，从不同角度衡量聚类质量。本题覆盖四种主要评估方法。

    \vspace{0.2cm}
    \textbf{问题：}
    \begin{enumerate}[(1)]
        \item \textbf{轮廓系数（Silhouette Coefficient）}：

        对样本 $i$，定义：
        \[
        a(i) = \text{样本 } i \text{ 到同簇内其他所有点的平均距离（簇内凝聚度）}
        \]
        \[
        b(i) = \text{样本 } i \text{ 到最近其他簇所有点的平均距离（簇间分离度）}
        \]
        \[
        s(i) = \frac{b(i) - a(i)}{\max(a(i), b(i))}
        \]
        整体轮廓系数为所有样本 $s(i)$ 的均值，取值范围为 $[-1, 1]$。

        请回答：
        \begin{enumerate}[a.]
            \item $s(i)$ 接近 $+1$、接近 $0$、接近 $-1$ 时，各自意味着什么？
            \item 轮廓系数\textbf{越大越好}还是越小越好？
        \end{enumerate}

        \item \textbf{簇内平方和（WCSS）与肘部法（Elbow Method）}：

        WCSS 定义为：
        \[
        \text{WCSS} = \sum_{k=1}^{K} \sum_{x_i \in C_k} \|x_i - \mu_k\|^2
        \]

        \begin{enumerate}[a.]
            \item 随着 $K$ 增大，WCSS 的变化趋势是什么？为什么 WCSS 不能直接用于选最优 $K$（即为什么不能直接最小化 WCSS）？
            \item 肘部法的思路是什么？"肘部"出现在哪里时，认为该 $K$ 值是合理选择？
        \end{enumerate}

        \item \textbf{CH 指数（Calinski-Harabasz Index）}：

        \[
        \text{CH}(K) = \frac{\text{BSS}/(K-1)}{\text{WSS}/(n-K)}
        \]
        其中 BSS 是簇间方差（Between-Cluster Sum of Squares），WSS 是簇内方差。CH 指数\textbf{越大越好}还是越小越好？它衡量的核心思想是什么？

        \item 用 \texttt{sklearn} 计算这三个指标的代码分别是什么？（只需写出关键一行调用）
    \end{enumerate}

    \vspace{0.2cm}
    {\color{gray}\small\textit{来源溯源：[文档第 10 章 10.1.4 聚类模型评估]}}
\end{problembox}

\begin{beginnerbox}[小白必读：为什么聚类评估比分类评估难？]
    \begin{itemize}
        \item \textbf{没有"标准答案"对比}：分类可以直接算"预测是否等于真实标签"，而聚类只有数据本身，没有外部标签告诉我们分几类是"对的"。
        \item \textbf{评估指标衡量的是"结构好不好"}：轮廓系数看"同类聚不聚、异类散不散"；WCSS 看簇内的紧凑程度；CH 指数看两者的比值。不同指标从不同角度打分，往往需要综合参考。
        \item \textbf{$K$ 的选取是一门"艺术"而非"科学"}：没有唯一正确答案，肘部法是经验性的，有时肘部不明显，需要结合业务知识判断。
    \end{itemize}
\end{beginnerbox}

\begin{solutionbox}[参考答案与详细解析]
    \textbf{(1) 轮廓系数解读：}

    \begin{itemize}
        \item \textbf{$s(i) \approx +1$}：$b(i) \gg a(i)$，即该点离自己所在簇的其他点很近（内聚好），离最近的其他簇很远（分离好）——分配得很好。
        \item \textbf{$s(i) \approx 0$}：$b(i) \approx a(i)$，该点处于两个簇的边界附近，分配到哪个簇差别不大。
        \item \textbf{$s(i) \approx -1$}：$a(i) \gg b(i)$，该点到自己簇内的距离比到最近其他簇的距离还大——说明该点被分配到了错误的簇。
    \end{itemize}

    \textbf{轮廓系数越大越好}，理想值接近 $+1$。

    \textbf{(2) WCSS 与肘部法：}

    \begin{itemize}
        \item \textbf{WCSS 随 $K$ 增大单调递减}：$K=n$（每个点单独一个簇）时，WCSS $= 0$，但这无任何意义。不能直接最小化 WCSS，否则答案永远是"每个点自己一个簇"。
        \item \textbf{肘部法}：绘制 $K$ 从 1 到若干取值时的 WCSS 曲线。随着 $K$ 增大，WCSS 下降速度先快后慢，形成"肘部"形状。在下降速度明显变缓的"转折点"处，该 $K$ 值被认为是数据中自然簇数的较好估计——继续增大 $K$ 带来的收益已经很小，不值得增加模型复杂度。
    \end{itemize}

    \textbf{(3) CH 指数：}

    \textbf{CH 指数越大越好}。其核心思想是衡量\textbf{簇间分散度与簇内紧凑度的比值}：分子（簇间方差大）意味着不同簇之间距离远，分母（簇内方差小）意味着同一簇内部紧凑——这正是好聚类的两个条件。

    \textbf{(4) sklearn 关键代码：}

    \begin{lstlisting}[language=Python]
from sklearn.metrics import silhouette_score, calinski_harabasz_score

# 轮廓系数（越大越好，范围 [-1, 1]）
sil = silhouette_score(X_scaled, labels)

# CH 指数（越大越好）
ch = calinski_harabasz_score(X_scaled, labels)

# WCSS（通过 KMeans 的 inertia_ 属性直接获取）
wcss = kmeans.inertia_
    \end{lstlisting}
\end{solutionbox}


%% ============================================================
\Needspace{5\baselineskip}
\begin{problembox}[Q10-06\quad 奇异值分解（SVD）：矩阵分解的核心原理]
    \textbf{题目描述：}\\
    奇异值分解（Singular Value Decomposition, SVD）是线性代数中最重要的矩阵分解方法之一，也是 PCA 降维的数学基础。理解 SVD 的结构是掌握降维的关键。

    \vspace{0.3cm}
    \textbf{题目背景：}\\
    对任意实矩阵 $A \in \mathbb{R}^{m \times n}$，SVD 将其分解为：
    \[
    A = U \Sigma V^T
    \]
    其中：
    \begin{itemize}
        \item $U \in \mathbb{R}^{m \times m}$：左奇异向量矩阵（列向量两两正交，$U^T U = I$）
        \item $\Sigma \in \mathbb{R}^{m \times n}$：对角矩阵，对角元素 $\sigma_1 \geq \sigma_2 \geq \cdots \geq \sigma_r \geq 0$ 称为奇异值
        \item $V \in \mathbb{R}^{n \times n}$：右奇异向量矩阵（列向量两两正交，$V^T V = I$）
    \end{itemize}

    \textbf{问题：}
    \begin{enumerate}[(1)]
        \item 矩阵 $A$ 的"秩"（Rank）与奇异值之间有什么关系？奇异值为 0 意味着什么？

        \item \textbf{截断 SVD 实现降维}：若只保留最大的 $k$ 个奇异值（$k < r$），得到近似矩阵：
        \[
        A_k = U_k \Sigma_k V_k^T
        \]
        其中 $U_k \in \mathbb{R}^{m \times k}$，$\Sigma_k \in \mathbb{R}^{k \times k}$，$V_k \in \mathbb{R}^{n \times k}$。

        \begin{enumerate}[a.]
            \item $A_k$ 的形状是什么？它是 $A$ 的精确复原还是近似复原？
            \item 为什么保留"最大的 $k$ 个奇异值"能保留最多信息？（奇异值大小代表什么？）
        \end{enumerate}

        \item 设 $A = \begin{pmatrix} 3 & 0 \\ 0 & 2 \\ 0 & 0 \end{pmatrix}$，直接写出其 SVD 分解结果（$U, \Sigma, V^T$），无需推导。（提示：$A$ 已经是接近对角形式）

        \item 用 NumPy 进行 SVD 分解的代码是什么？如何只保留前 $k=2$ 个奇异值来实现近似重建？
    \end{enumerate}

    \vspace{0.2cm}
    {\color{gray}\small\textit{来源溯源：[文档第 10 章 10.2.1 奇异值分解]}}
\end{problembox}

\begin{metaphorbox}[生活比喻：SVD 就像"拍照压缩"]
    你拍了一张 $1000 \times 1000$ 的图片（100 万个像素值的矩阵 $A$）：

    \begin{itemize}
        \item \textbf{完整 SVD}：把图片"拆解"成一系列"基础图案"的叠加，每个基础图案由一对奇异向量定义，叠加的"权重"就是奇异值。
        \item \textbf{奇异值大 = 基础图案重要}：奇异值大的那几个"基础图案"贡献了图片的主要结构（轮廓、主色调）；奇异值小的贡献的是细节和噪声。
        \item \textbf{截断 SVD = 有损压缩}：只保留最重要的 $k$ 个基础图案，用 $A_k \approx A$ 来存储图片。$k$ 越小，文件越小，但图片越模糊；$k$ 越大，越接近原图，但存储空间也越大。
    \end{itemize}
\end{metaphorbox}

\begin{solutionbox}[参考答案与详细解析]
    \textbf{(1) 秩与奇异值的关系：}

    矩阵 $A$ 的秩 $r$ 等于非零奇异值的个数：$r = \text{rank}(A) = |\{i : \sigma_i > 0\}|$。奇异值为 0 意味着该方向上矩阵没有"信息量"（数据在该方向上的投影为零），对应的奇异向量对矩阵的重建没有贡献。

    \textbf{(2) 截断 SVD：}

    \begin{enumerate}[a.]
        \item $A_k = U_k \Sigma_k V_k^T$ 的形状仍为 $\mathbb{R}^{m \times n}$（与原矩阵相同）。它是\textbf{近似复原}而非精确复原——用更少的参数（$k$ 个奇异值及对应向量）近似表示原矩阵，$k < r$ 时会有信息损失。
        \item 奇异值 $\sigma_i$ 的大小代表矩阵在该奇异向量方向上的"能量"（方差贡献）。最大的奇异值对应最能解释数据变化的方向，保留大奇异值意味着保留了原矩阵中最主要的结构信息，舍弃小奇异值相当于去除噪声或次要细节。
    \end{enumerate}

    \textbf{(3) 简单矩阵的 SVD：}

    $A = \begin{pmatrix} 3 & 0 \\ 0 & 2 \\ 0 & 0 \end{pmatrix}$ 已为列正交形式，其 SVD 可直接写出：

    \[
    U = \begin{pmatrix} 1 & 0 & 0 \\ 0 & 1 & 0 \\ 0 & 0 & 1 \end{pmatrix},\quad
    \Sigma = \begin{pmatrix} 3 & 0 \\ 0 & 2 \\ 0 & 0 \end{pmatrix},\quad
    V^T = \begin{pmatrix} 1 & 0 \\ 0 & 1 \end{pmatrix}
    \]

    奇异值为 $\sigma_1 = 3$，$\sigma_2 = 2$，秩 $r = 2$。

    \textbf{(4) NumPy SVD 代码：}

    \begin{lstlisting}[language=Python]
import numpy as np

A = np.array([[3, 0], [0, 2], [0, 0]], dtype=float)

# 完整 SVD 分解
U, s, Vt = np.linalg.svd(A, full_matrices=True)
# U: (3,3)  s: (2,) 奇异值向量  Vt: (2,2)

# 只保留前 k=2 个奇异值进行近似重建
k = 2
U_k = U[:, :k]         # (3, 2)
S_k = np.diag(s[:k])   # (2, 2)
Vt_k = Vt[:k, :]       # (2, 2)

A_k = U_k @ S_k @ Vt_k   # 近似重建，形状仍为 (3, 2)
print("重建误差:", np.linalg.norm(A - A_k))
    \end{lstlisting}
\end{solutionbox}


%% ============================================================
\Needspace{5\baselineskip}
\begin{problembox}[Q10-07\quad 主成分分析（PCA）：用 SVD 实现降维的完整流程]
    \textbf{题目描述：}\\
    主成分分析（Principal Component Analysis, PCA）是最常用的降维方法，其核心可以用 SVD 来实现。本题通过一个完整的数学推导和代码实践，理解 PCA 的本质。

    \vspace{0.3cm}
    \textbf{题目背景：}\\
    设数据矩阵 $X \in \mathbb{R}^{n \times p}$（$n$ 个样本，$p$ 个特征）。PCA 的目标是找到一组新的低维坐标轴（主成分），使数据在这些坐标轴上的\textbf{方差最大化}（信息保留最多）。

    SVD 实现 PCA 的步骤：
    \begin{enumerate}
        \item \textbf{中心化}：$X_c = X - \bar{X}$（每列减去对应特征的均值）
        \item \textbf{对 $X_c$ 做 SVD}：$X_c = U \Sigma V^T$
        \item \textbf{主成分方向}：$V$ 的列向量（即 $V^T$ 的行向量）就是主成分方向
        \item \textbf{降维投影}：$Z = X_c V_k$（保留前 $k$ 个主成分，$Z \in \mathbb{R}^{n \times k}$）
    \end{enumerate}

    \textbf{问题：}
    \begin{enumerate}[(1)]
        \item 为什么 PCA 第一步必须\textbf{中心化}（均值归零）？不中心化会发生什么？

        \item "方差最大化"和"信息保留最多"为什么是同一回事？请用直觉解释（不需要数学推导）。

        \item 在 PCA 中，如何判断降到 $k$ 维后\textbf{保留了多少原始信息}？（用奇异值表示）给出公式。

        \item 下面是用 SVD 实现 PCA 并降维鸢尾花数据的代码，请填写空白处：

        \begin{lstlisting}[language=Python]
import numpy as np
from sklearn.datasets import load_iris

iris = load_iris()
X = iris.data   # shape: (150, 4)

# Step 1: 中心化
X_c = X - X.mean(axis=____)   # 填写 axis 参数

# Step 2: SVD 分解
U, s, Vt = np.linalg.svd(X_c, full_matrices=False)

# Step 3: 降到 2 维（保留前 2 个主成分）
k = 2
Z = X_c @ Vt[:k, :].____   # 填写转置操作

# Step 4: 计算前 2 个主成分解释的方差比例
explained_ratio = (s[:k]**2).sum() / ____   # 填写分母

print(f"降维后形状: {Z.shape}")
print(f"前2个主成分解释方差比例: {explained_ratio:.2%}")
        \end{lstlisting}

        \item PCA 中的\textbf{主成分}（Principal Component）在几何上代表什么？第一主成分方向和数据方差最大方向有什么关系？
    \end{enumerate}

    \vspace{0.2cm}
    {\color{gray}\small\textit{来源溯源：[文档第 10 章 10.2.2 主成分分析]}}
\end{problembox}

\begin{beginnerbox}[小白必读：PCA 到底在做什么？]
    \begin{itemize}
        \item \textbf{本质是"换一套坐标系"}：原来的坐标轴是人为定义的（身高、体重、年龄……），主成分轴是数据自己告诉我们的——沿着数据变化最大的方向。
        \item \textbf{投影 = 在新坐标系下读坐标}：$Z = X_c V_k$ 的含义是：把每个数据点投影到前 $k$ 个主成分方向上，得到新的低维坐标。
        \item \textbf{为什么要降维？}：4 维的鸢尾花数据很难直接可视化，降到 2 维后，散点图上就能清晰看到三个品种的分布，发现聚类结构。高维特征中往往存在大量冗余（相关的特征），PCA 能把这些冗余压缩掉，只保留最有信息量的方向。
    \end{itemize}
\end{beginnerbox}

\begin{metaphorbox}[生活比喻：PCA 就像"找最佳拍照角度"]
    你有一个 3D 雕塑，想拍一张 2D 照片，让看照片的人尽可能理解雕塑的形状：

    \begin{itemize}
        \item \textbf{随意拍（随机降维）}：可能正对着雕塑最薄的一面，看起来像一条线，信息丢失严重。
        \item \textbf{PCA 拍法}：自动找到能让雕塑"展开得最大"的角度——从这个角度看，雕塑的细节、轮廓、层次感保留最多。
        \item \textbf{第一主成分}：信息最丰富的那个维度（雕塑展开最大的方向）；第二主成分：在与第一主成分垂直的方向中，信息最丰富的那个；以此类推。
    \end{itemize}

    \textbf{PCA 的精髓}：用更少的维度，保留最多的"信息量"（方差）。
\end{metaphorbox}

\begin{solutionbox}[参考答案与详细解析]
    \textbf{(1) 中心化的必要性：}

    若不中心化，SVD 的结果会"偏向"数据均值的方向——第一个奇异向量会主要指向数据的均值方向，而非真正的最大方差方向。中心化的目的是消除均值的影响，让 SVD 纯粹分析数据的\textbf{离散/变化结构}，而非数据的位置。

    \textbf{(2) 方差最大化 = 信息保留最多：}

    如果数据在某个方向上的投影方差很大，说明不同样本在这个方向上有显著差异——这个差异就是可以用来区分样本的信息。如果投影方差接近零，所有样本投影到这个方向上几乎得到同一个值，这个方向传达不了任何样本间的区别，相当于没有信息。因此，保留方差最大的方向 = 保留最多能区分样本的信息。

    \textbf{(3) 方差解释比例公式：}

    前 $k$ 个主成分解释的方差比例为：
    \[
    \text{Explained Variance Ratio} = \frac{\sum_{i=1}^{k} \sigma_i^2}{\sum_{i=1}^{r} \sigma_i^2}
    \]
    其中 $\sigma_i$ 是第 $i$ 个奇异值，$r$ 是非零奇异值的个数（即矩阵的秩）。奇异值的平方 $\sigma_i^2$ 正比于数据在第 $i$ 个主成分方向上的方差。

    \textbf{(4) 代码填空答案：}

    \begin{lstlisting}[language=Python]
# Step 1: 中心化
X_c = X - X.mean(axis=0)   # axis=0 表示对每列（每个特征）求均值

# Step 3: 降到 2 维
Z = X_c @ Vt[:k, :].T      # Vt[:k,:].T 形状为 (4, 2)，投影到 2 维

# Step 4: 方差解释比例
explained_ratio = (s[:k]**2).sum() / (s**2).sum()   # 分母是全部奇异值的平方和
    \end{lstlisting}

    运行结果说明：鸢尾花 4 维数据降到 2 维后，通常前 2 个主成分可以解释约 97\% 的方差，说明降维几乎没有丢失有用信息。

    \textbf{(5) 主成分的几何含义：}

    \begin{keybox}[主成分的几何解释]
        \begin{itemize}
            \item \textbf{主成分方向}：$V$ 的列向量，彼此正交（垂直）的单位向量，构成新的坐标轴。
            \item \textbf{第一主成分（PC1）}：数据点在该方向上的投影方差最大，即这个方向是数据"变化最剧烈"的方向。
            \item \textbf{第二主成分（PC2）}：在与 PC1 垂直的所有方向中，投影方差最大的方向。
            \item \textbf{降维的本质}：把原来 $p$ 维坐标系旋转为以主成分为轴的新坐标系，然后只保留前 $k$ 个轴的坐标，丢弃方差最小（信息最少）的若干维度。
        \end{itemize}
    \end{keybox}
\end{solutionbox}

\begin{appbox}[现实应用场景]
    \textbf{PCA 和聚类的经典组合——人脸识别、基因表达分析：}\\
    在人脸识别中，原始图片可能是 $128 \times 128 = 16384$ 维的向量，直接聚类或分类效率极低。先用 PCA 降到 100 维，再做 K-means 聚类或分类，可以在大幅降低计算量的同时保留 99\% 以上的信息。这种"先 PCA 降维、再聚类/分类"的两阶段流程，是工业界数据分析的经典 pipeline。SVD 在推荐系统中也扮演核心角色——Netflix、Amazon 的协同过滤算法，本质上就是对用户-商品评分矩阵进行截断 SVD，从中挖掘用户偏好的隐含特征。
\end{appbox}


% ===== 本章小结 =====
\section{本章小结：无监督学习的两条主线}

学完本章 7 道题，可以把无监督学习的核心浓缩成两条主线：

\begin{itemize}
    \item \textbf{聚类主线}：数据有内在分组结构 $\to$ K-means（球形簇，需要预设 $K$）/ 层次聚类（树状结构，可视化好）/ DBSCAN（任意形状，能识别噪声）$\to$ 用轮廓系数、WCSS、CH 指数评估聚类质量 $\to$ 肘部法辅助选 $K$。
    \item \textbf{降维主线}：高维数据冗余多 $\to$ SVD 将矩阵分解为奇异值和奇异向量 $\to$ 截断 SVD 只保留大奇异值方向 $\to$ PCA 通过对中心化数据做 SVD 找到方差最大方向 $\to$ 降维后保留最多信息。
\end{itemize}

\begin{metaphorbox}[给零基础读者的一句总比喻]
    \textbf{如果把数据比作一群人在广场上的站位：}
    \begin{itemize}
        \item \textbf{聚类}：不给任何提示，让你把自然地聚在一起的人归为一组——K-means 用圆圈圈，DBSCAN 用"密度围栏"，层次聚类画一棵家谱树。
        \item \textbf{PCA 降维}：广场是 3D 的，但你只有一张 2D 纸，选一个最能体现人群分布的角度拍下来——PCA 帮你找到这个最佳角度，让拍出来的平面图损失最少信息。
        \item \textbf{两者结合}：先 PCA 降维（消除噪声、减少冗余维度），再聚类（在低维空间中结构更清晰），是无监督分析的黄金搭档。
    \end{itemize}
\end{metaphorbox}

\end{document}
